%% Section 3: Setup and notations for the height inequality

In the next few sections we will prove Theorem~\ref{ThmHtInequality}.
Let us first fix  the setting.

All varieties are over an algebraically closed subfield $k$ of $\IC$.
The ambient data is given as above Theorem~\ref{ThmHtInequality}. We
repeat it here.
%The ambient data, with $\cL$ as in Raynaud's~\cite{LNM119}, is given as follows:
\begin{itemize}
  \item Let $S$ be a regular, irreducible, quasi-projective variety over  $k$
  that is Zariski open in an irreducible projective variety
  $\overline{S}\subset \IP_k^m$.
%  \Pinlinecomment{$\overline{S}$ used to be regular, is it safe to    assume that it may have horrible singularities?} \Gcomment{I don't think we need $\bar{S}$ to be regular in the proof. In fact we don't really work with the boundaries. And when applying the height inequality in Prop 7.1, we cannot assume $\bar{S}$ to be regular because $S$ is a subvariety of $\mathbb{M}_g$, and its Zariski closure in $\overline{\mathbb{M}_g}$ may not be regular even if $\overline{\mathbb{M}_g}$ is.}
  % \Pcomment{Question to self: Revert to $S$ quasi-projective?} \Gcomment{Affine changed to quasi-projective.}%\Gcomment{Seems better to me. The immersion in the second bullet point can be obtained from $\cL \otimes \pi^*\cM$ with $\cL$ as in Remark~\ref{rmk:projembedding} and some $\cM$ very ample on $S$; see Stack Project.}
  \item Let $\pi\colon \cA\rightarrow S$ be an abelian scheme
     presented by a closed immersion $\cA\rightarrow\IP_k^n\times S$ over
     $S$. %!!!FURTHER DOWN WE SOMETIMES HAVE: $S\times \IP_k^n$. !!!FIXME
  \item
    From the previous point, we get a closed immersion of the generic
    fiber $A$ of $\cA\rightarrow S$  into  $\IP_{k(S)}^n$.
    We assume that $A\rightarrow \IP_{k(S)}^n$  
    arises from a basis of the global sections of 
    the $l$-th power $L$ of a symmetric ample
    line bundle with $l\ge 4$. 
  \item
    Finally, we assume \texttt{(Hyp)} as on page \pageref{def:hyp}.
\end{itemize}
% of  relative dimension $g$, [PH]: relative dimension never used
%with $0$-section $\epsilon\colon S\rightarrow\cA$, and [PH]:
%$0$-section is never used
%% \item let $L$ be a line bundle on $\cA_\eta$ that is the fourth power
%% of an ample line bundle on $\cA_\eta$ and $[-1]^* L \cong L$. 
%% \item let $\cL$ be a line bundle on $\cA$ that is $S$-ample and such
%% that $[-1]^*\cL\cong \cL$. 
%such that  $\epsilon^*\cL$ is trivial, and such that
%$[-1]^*\cL \cong \cL$.

%If $S$ is a regular, irreducible, quasi-projective variety of $k$,
%then $\overline{S}$ as in the first bullet point always exists due to
%Hironaka's Theorem.\Pinlinecomment{Hironaka is not necessary if we
%  don't ask $\overline{S}$ to be regular. But this assumption may be important...}
From the third bullet power, we see that the image of $A$ is
projectively normal in $\IP_{k(S)}^n$, \textit{cf}.~\cite[Theorem 9]{Mumford:quadeq}. By the fourth bullet point,
Proposition~\ref{PropBettiForm} provides the Betti form $\omega$ on
$\an{\cA}$.


For $s\in S(k)$ we write $\cA_s$ for the abelian variety $\pi^{-1}(s)$. 

\begin{remark}
\label{rmk:projembedding}
Let $S$ be as in the first bullet point. %For simplicity we assume in addition that $S$ is affine, this suffices for Theorem~\ref{ThmHtInequality}. 
Let
$\pi \colon \cA \rightarrow S$ be an abelian scheme.
Suppose $L_0$ is a 
symmetric and ample line bundle on $A$, the generic fiber of $\pi$.
An immersion of $\cA$ as in the second bullet point can  be obtained
as follows. 
 By \cite[Th\'eor\`eme~XI~1.13]{LNM119}  there exists an
$S$-ample line bundle $\cL$ on $\cA$ whose restriction to the
generic fiber of $\cA\rightarrow S$ is isomorphic to $L_0^{\otimes l}$
for some integer $l\ge 4$. 
We may assume in addition that $\cL$
satisfies
$[-1]^* \cL \cong \cL$ and even that $\cL$ becomes trivial when pulled back under the
zero section $S\rightarrow \cA$, see \cite[Remarque~XI~1.3a]{LNM119}.
%We may even assume that $\cL$ 
After replacing $\cL$
by  a sufficiently high power, we may assume that $\cL$ is very ample over
$S$.
We fix a basis of global sections of $L_0^{\otimes l}$ and, as $l\ge
4$, thereby
realize the
 generic fiber of $\pi$ as a projectively normal subvariety of 
$\IP_{k(S)}^n$.
Now we can take $L$ in the third bullet point to be $L_0^{\otimes l}$,
which is the restriction of $\cL$ to $A$.
A closed immersion $\cA \rightarrow \IP^n_k \times S$ as in the second
bullet point %exists because $S$ is affine. 
 arises from 
$\cL \otimes \pi^*\cM$ for some very ample line bundle $\cM$ on $S$; see \cite[Proposition~4.4.10.(ii) and Proposition~4.1.4]{EGAII}.
On restricting to a fiber of $\cA\rightarrow S$ the induced closed immersion  $\cA_s\rightarrow\IP^n_k$ comes from the restriction $\cL|_{\cA_s}$. 
\end{remark}


Write $\overline \cA$ for the Zariski closure
of $\cA$ in $\IP_k^n\times \overline{S}$. Then $\overline \cA$ is irreducible
but not necessarily regular.
On any product of $r$ projective spaces and if $a_1,\ldots,a_r\in\IZ$,
we let  $\cO(a_1,\ldots,a_r)$ denote the tensor product over all $i\in
\{1,\ldots, r\}$ of the
pull-back under the $i$-th projection of $\cO(a_i)$. 
We write $\cL$ for the restriction of $\cO(1,1)$ to $\overline{\cA}$. 

\subsection{Height functions on $\cA$}\label{SubsectionHeight}
If $k=\IQbar$ we have several height functions on $\cA(\IQbar)$.

For any $n \in \IN$, we always consider the absolute logarithmic 
Weil height function $\IP_{\IQbar}^n(\IQbar) \rightarrow \IR$, or just 
Weil height, defined as in \cite[$\mathsection$1.5.1]{BG}.

Now say $P\in \cA(\IQbar)$, we write $P=(P',\pi(P))$ with
$P'\in \IP_{\IQbar}^n(\IQbar)$ and
$\pi(P)\in \IP_{\IQbar}^m(\IQbar)$. The sum of
Weil heights
\begin{equation}\label{EqHeightTotal}
h(P) = h(P') + h(\pi(P))
\end{equation}
defines our first height $\cA(\IQbar)\rightarrow [0,\infty)$ which we
call the \textit{naive height} on $\cA$. It depends on the fixed
immersion of $\cA$.

The line bundle $[-1]^*\cL|_{\cA} \otimes\cL|_{\cA}^{\otimes -1}$ of
${\cA}$ restricted to the generic fiber $A$ of $\cA\rightarrow S$
equals $[-1]^*L \otimes L^{\otimes -1}$. By the third bullet point above this line
bundle is trivial. So it equals $\pi^*\cK$ for some line bundle $\cK$
of $S$ by \cite[Corollaire~21.4.13 (pp.~361 of EGA IV-4, in Errata et Addenda, liste~3)]{EGAIV}. We conclude
$[-1]^*\cL|_{\cA_s}\cong \cL|_{\cA_s}$ for all $s\in S(\IQbar)$. So
the function (\ref{EqHeightTotal}) represents the height function,
defined up-to $O_s(1)$, given by the Height
Machine, \textit{cf.} \cite[Theorem 2.3.8]{BG}, applied to
$(\cA_s,\cL_s)$. As $\cL_s$ is symmetric, the \textit{fiberwise
N\'eron--Tate} or \textit{canonical height} $\hat
h_{\cA} \colon \cA(\IQbar)\rightarrow [0,\infty)$, defined by the
convergent limit
\begin{equation}
\label{eq:fiberwiseNT}
\hat h_{\cA}(P) = \lim_{N\rightarrow\infty} \frac{h([N](P))}{N^2},%\hat h_{\cA_{\pi(P)}} (P).
\end{equation}
is a quadratic form on $\cA_s(\IQbar)$. In the notation \cite[Chapter~9]{BG} the height
(\ref{eq:fiberwiseNT}) is $\hat h_{\cA_s,\cL_s}$ where $s=\pi(P)$. 


\begin{remark}\label{RemarkHtcompatible}
  We use here the notation of Remark~\ref{rmk:projembedding}.
  So that the immersion $\cA\rightarrow \IP^n_S$ arises via
  $L_0^{\otimes l}$. To normalize,
  we divide (\ref{eq:fiberwiseNT}) by $l$ and obtain the 
  N\'eron--Tate height $\hat h_{\cA,L_0}\colon \cA(\IQbar)\rightarrow
  [0,\infty)$.

  Let us verify that $\hat h_{\cA,L_0}$ depends only on $L_0$. Suppose
  $\cL'$ is another line bundle on $\cA$ that restricts to
  $L_0^{\otimes l}$,
  then $\cL'\otimes\cL^{\otimes -1}$ is trivial on $A$. By \cite[Corollaire~21.4.13 (pp.~361 of EGA IV-4, in Errata et Addenda, liste~3)]{EGAIV}, this difference is the pull-back of some line
  bundle on $S$ under $\cA\rightarrow S$. So the restriction of
  $\cL'\otimes\cL^{\otimes -1}$ to $\cA_s$ for each $s\in S(\IC)$ is
  trivial. Thus $\cL|_{\cA_s}$ and $\cL'|_{\cA_s}$ induce the same N\'eron--Tate
  height on $\cA_s(\IQbar)$, see~\cite[$\mathsection$9.2]{BG}. 
\end{remark}

\subsection{Integration against the Betti form}\label{SubsectionSetUpHtIneqBetti}
Let $\cA$  and $S$ be as in the beginning of this section, so they are
defined over an algebrically closed subfield $k$ of $\IC$. 
Recall that $\omega$ is the Betti form on $\an{\cA}$ as provided by
Proposition~\ref{PropBettiForm}. In particular, it is a
semi-positive $(1,1)$-form on $\an{\cA}$ such that $[N]^*\omega =
N^2 \omega$ for all $N\in\IZ$. We discuss here a modification of  the
Betti form that has compact support.

Fix $X$ to be an irreducible closed subvariety of $\cA$ of dimension $d$, such that $\pi|_X \colon X \rightarrow S$ is dominant. 

We are not allowed to integrate $\omega^{\wedge d}$ over $\sman{X}$ as $\omega^{\wedge d}$ may not have compact support. So we modify $\omega$ in the following way.

Suppose we are provided with a base point $s_0 \in \an{S}$. Let
furthermore $\Delta$ be a relatively compact, contractible,
open neighborhood of $s_0$ in $\an{S}$. Denote by $\cA_{\Delta}$ the open
subset $\pi^{-1}(\Delta)$ of $\an{\cA}$. Fix a smooth bump function $\vartheta \colon \an{S}\rightarrow [0,1]$ with compact support $K \subset \Delta$ such that $\vartheta(s_0)=1$. Finally, we define $\theta = \vartheta\circ\pi \colon \an{\cA} \rightarrow [0,1]$.
Then $\theta\omega$ is a semi-positive smooth $(1,1)$-form on $\an{\cA}$;
unlike the Betti form, it may not be closed. By construction, the support
of $\theta\omega$ lies in $\pi^{-1}(K)$ which is compact as $\pi$ is proper and $K$ is compact. 

\begin{remark}
  \label{rmk:nondeg}
  Suppose $X$ is non-degenerate, namely $X$ satisfies one of the two
  equivalent conditions in property (iii) of
  Proposition~\ref{PropBettiForm}. Then $\an{X}$ contains a smooth
  point $P_0$ at which $\omega|_{\sman{X}}^{\wedge d}>0$. Then we will
  take $s_0=\pi(P_0)$.
\end{remark}



\subsection{The graph construction}
\label{ss:graph}

Let $N\in\IZ$. The multiplication-by-$N$ morphism $[N] \colon \cA\rightarrow \cA$ may not
extend to a morphism $\overline{\cA}\rightarrow \overline{\cA}$. We
overcome this by using the graph construction.

Recall that we have identified
$\cA\subset\overline{\cA}\subset \IP_k^n\times \overline{S}
\subset\IP_k^n\times\IP_k^m$. 

We write
$\rho_1,\rho_2\colon \IP_k^n\times\IP_k^n \times \IP_k^m \rightarrow
\IP_k^n\times \IP_k^m$ for the
 two projections $\rho_1(P,Q,s)=(P,s)$ and $\rho_2(P,Q,s)=(Q,s)$.


Consider $\Gamma_N$ the graph of $[N]$, determined by 
\begin{equation*}
  \Gamma_N = \{ (P,[N](P)) : P \in \cA(k) \}. 
\end{equation*}
We consider it as an irreducible closed subvariety of
$\cA\times_S\cA$. 

Let $X$ be an irreducible closed subvariety of $\cA$ of dimension $d$. 
The graph $X_N$ of $[N]$ restricted to $X$ is
an irreducible closed  subvariety of $\Gamma_N$ determined by 
\begin{equation*}
   \{(P,[N]P): P \in X(k) \}.
\end{equation*}


Observe that $\rho_1|_{\Gamma_N} \colon \Gamma_N\rightarrow \cA$ is an isomorphism; it maps $(P,[N](P))$ to $P$.  % In particular,  $\Gamma_N$ is a complex manifold.
So we can use $\rho_1|_{\Gamma_N}^{-1}$ to identify $X$ with $X_N$.


Moreover, $\rho_2|_{\Gamma_N}$ maps $(P,[N](P))$ to $[N](P)$. Therefore
\begin{equation}
\label{eq:relatetoN}
  \rho_2|_{\Gamma_N} \circ \rho_1|_{\Gamma_N}^{-1} = [N].
\end{equation}


Let $\overline{X_N}$ be the Zariski closure of $X_N$ in
$\overline{\cA}\times_{\overline{S}}\overline{\cA}\subset
\IP^n_{\overline{S}}\times_{\overline{S}}\IP^n_{\overline{S}} =
\IP_k^n\times\IP_k^n\times \overline{S}
\subset\IP_k^n\times\IP_k^n\times \IP_k^m$.  Then $\overline{X_N}$ is an irreducible projective variety (which is not necessarily regular) with $\dim \overline{X_N} = \dim X_N = \dim X$.

% Let $\rho_0 \colon  \IP_k^n\times\IP_k^n\times \IP_k^m  \rightarrow \IP_k^m$ be the projection to the final factor.
In the next section, we will use the following line bundles on $\overline{X_N}$. Define
\begin{equation}\label{def:cF}
\cF = \rho_2^*\cO(1,1)|_{\overline{X_N}} = \cO(0,1,1)|_{\overline{X_N}}
\end{equation}
and
\begin{equation}\label{def:cM}
  \cM %= \rho_0^*\cO(1)|_{\overline{X_N}}
  = \cO(0,0,1)|_{\overline{X_N}}.
\end{equation}

Let us close this subsection by relating the height functions defined
by $\cF$ and $\cM$ with the ones in
$\mathsection$\ref{SubsectionHeight}. Assume $k = \IQbar$. Let $P \in
X(\IQbar)$. Write $P = (P',\pi(P))$ with
 $P' \in \IP_\IQbar^n(\IQbar)$ and $\pi(P) \in\IP_\IQbar^m(\IQbar)$. %  for $\cA
% \subseteq  \IP_\IQbar^n\times S \subseteq \IP_\IQbar^n \times \IP_\IQbar^m$.
We have $[N](P) = ( P'_N,\pi(P))$ for some $P'_N \in \IP_\IQbar^n(\IQbar)$.

Under the immersion $\overline{X_N} \subseteq \IP_\IQbar^n \times
\IP_\IQbar^n \times \IP_\IQbar^m$, the point $(P,[N]P)$ in
$\overline{X_N}(\IQbar)$ becomes $P_N=(P',P'_N,\pi(P)) \in
(\IP_\IQbar^n \times \IP_\IQbar^n \times \IP_\IQbar^m)(\IQbar)$.
The function $P_N\mapsto h([N](P)) = h(P'_N)+h(\pi(P))$ defined in
\eqref{EqHeightTotal}
represents the height attached by the Height Machine to
$(\overline{X_N},\cF)$
and $P_N\mapsto h(\pi(P))$ represents the height attached to $(\overline{X_N},\cM)$.

%%% Local Variables:
%%% TeX-master: "main"
%%% End:

